\section{Appendix Figures}
\label{app:figures}

\renewcommand{\thefigure}{D.\arabic{figure}}
\setcounter{figure}{0}


\begin{figure}[h]
    \begin{center}
    \caption{Illustration of the Trade-off between Worker Specialization and Coordination in Task Assignment}
    \label{fig:intro_task_illustration}
    \includegraphics[width=1\linewidth]{plots/job_design.png}
    \end{center}
    \footnotesize
    \emph{Notes:} The blue bounded rectangles indicate tasks.
    The height of each rectangle corresponds to the task's skill requirement ($\skillcost{}$) and its width to the time requirement ($\timecost{}$).
    The shaded areas represent the wage bills of jobs.
    In the left panel, tasks are assigned to two specialized workers but a hand-off cost (the pink rectangle) is introduced due to coordination frictions between them.
    In the right panel, tasks are bundled into a single job performed by one worker, eliminating hand-off costs but requiring a worker skilled enough to perform both tasks and compensated with a higher per-unit-time wage rate.
\end{figure}




\newpage
\begin{figure}[ht!]
    \caption{Task Sequence of Computer Programmers Occupation}
    \label{fig:computer_programmers}
    \begin{center}
    \resizebox{\textwidth}{!}{
    \includegraphics[width=\linewidth]{plots/taskSequence_vs_anthropicIndex/all_occupation_task_sequences/task_sequence_15-1251.00_Computer_Programmers.png}
    }
    \end{center}
    \footnotesize{
    \emph{Notes:} This figure shows the task sequence for an occupation with a high share of its tasks (about two thirds) executed by AI. 
    The task ordering is generated by GPT-5-mini.
    The execution labels come from the Anthropic Economic Index. 
    The O*NET code for this occupation is 15-1251.
    }
\end{figure}

\newpage
\begin{figure}[ht!]
    \caption{Task Sequence of Public Relations Specialists Occupation}
    \label{fig:public_relations}
    \begin{center}
    \resizebox{\textwidth}{!}{
    \includegraphics[width=\linewidth]{plots/taskSequence_vs_anthropicIndex/all_occupation_task_sequences/task_sequence_27-3031.00_Public_Relations_Specialists.png}
    }
    \end{center}
    \footnotesize{
    \emph{Notes:} This figure shows the task sequence for an occupation with a moderate share of its tasks (about one third) executed by AI. 
    The task ordering is generated by GPT-5-mini. 
    The execution labels come from the Anthropic Economic Index. 
    The O*NET code for this occupation is 27-3031.
    }
\end{figure}

\newpage
\begin{figure}[ht!]
    \caption{Task Sequence of Electronic Equipment Installers and Repairers, Motor Vehicles}
    \label{fig:electronic_equipments}
    \begin{center}
    \resizebox{\textwidth}{!}{
    \includegraphics[width=\linewidth]{plots/taskSequence_vs_anthropicIndex/all_occupation_task_sequences/task_sequence_49-2096.00_Electronic_Equipment_Installers_and_Repairers__Motor_Vehicles.png}
    }
    \end{center}
    \footnotesize{
    \emph{Notes:} This figure shows the task sequence for an occupation with none of its tasks executed by AI. 
    The task ordering is generated by GPT-5-mini. 
    The execution labels come from the Anthropic Economic Index. 
    The O*NET code for this occupation is 49-2096.
    }
\end{figure}




\newpage
\begin{figure}[t!]
\caption{Relationships Between Occupational AI Exposure, Empirical Fragmentation, and AI Execution (Execution-Based Empirical Fragmentation Measure)}
\label{fig:fragmentation_index_regression_execution}

\begin{center}
\begin{subfigure}[b]{\textwidth}
    \centering
    \begin{subfigure}[b]{0.49\textwidth}
        \centering
        \captionsetup{labelformat=empty}
        \caption{(a) Empirical Fragmentation Index vs.\ AI Exposure}
        \includegraphics[width=\textwidth]{plots/fragmentationIndex/fragmentation_vs_aiExposure_execution.png}
    \end{subfigure}
    \hfill
    \begin{subfigure}[b]{0.49\textwidth}
        \centering
        \captionsetup{labelformat=empty}
        \caption{(b) AI Execution vs.\ AI Exposure}
        \includegraphics[width=\textwidth]{plots/fragmentationIndex/aiFraction_vs_aiExposure_execution.png}
    \end{subfigure}
\end{subfigure}

\end{center}
\footnotesize{
\emph{Notes:} These graphs show how the fragmentation of AI-able tasks within an occupation's workflow shapes the conversion of AI-exposed tasks into AI-executed tasks.  
Each bin represents a single O*NET occupation in both panels.  
Panel~(a) plots version~4 of the empirical fragmentation index defined in Table~\ref{tab:fragmentation_index_regression_execution} against the share of AI-exposed tasks, whereas Panel~(b) plots the share of realized AI-executed tasks against the share of AI-exposed tasks. 
All six highlighted occupations have 23 tasks in their production sequence and share similar occupation-level exposure to AI: \textit{Emergency Management Directors, Law Teachers,} and \textit{Philosophy and Religion Teachers} each have 39\% of their tasks exposed to AI, whereas \textit{Insurance Claims and Policy Processing
Clerks, Social Science Research Assistants,} and \textit{Communications Teachers} each have 35\% of their tasks exposed to AI.
}
\end{figure}




\newpage
\begin{figure}[ht!]
  \begin{center}
  \caption{Effect of Neighboring Tasks' AI Execution Status on Task's AI Execution Likelihood (GPT-filtered Sample)} 
  \label{fig:DWA_regression_aiExecution_GPTsample}
  \begin{subfigure}[b]{\textwidth}
    \captionsetup{labelformat=empty}
    \caption{Panel (A): No Fixed Effects}
    \includegraphics[width=\textwidth]{plots/execTypeVaryingDWA_noTasksWithRepetitiveDWAs/is_ai/AME_filtered_is_ai_no_fe_no_dwa.png}
  \end{subfigure}
  
  \vspace{1em}
  
  \begin{subfigure}[b]{\textwidth}
    \captionsetup{labelformat=empty}
    \caption{Panel (B): SOC Major Group Fixed Effects}
    \includegraphics[width=\textwidth]{plots/execTypeVaryingDWA_noTasksWithRepetitiveDWAs/is_ai/AME_filtered_is_ai_major_fe_no_dwa.png}
  \end{subfigure}
  
  \vspace{1em}
  
  \begin{subfigure}[b]{\textwidth}
    \captionsetup{labelformat=empty}
    \caption{Panel (C): SOC Minor Group Fixed Effects}
    \includegraphics[width=\textwidth]{plots/execTypeVaryingDWA_noTasksWithRepetitiveDWAs/is_ai/AME_filtered_is_ai_minor_fe_no_dwa.png}
  \end{subfigure}

  \vspace{1em}

  \begin{subfigure}[b]{\textwidth}
    \captionsetup{labelformat=empty}
    \caption{Panel (D): Detailed Work Activity Fixed Effects}
    \includegraphics[width=\textwidth]{plots/execTypeVaryingDWA_noTasksWithRepetitiveDWAs/is_ai/AME_filtered_is_ai_no_fe_with_dwa.png}
  \end{subfigure}
  \end{center}
  \footnotesize{\emph{Notes:} These graphs show that, among similar tasks appearing in multiple occupations, those whose immediate neighbors are AI-executed exhibit higher probabilities of being executed by AI, while more distant neighbors have little to no effect on a task's AI execution likelihood.
  The red dashed line in each graph indicates the observed average marginal effect of the corresponding variable in the original dataset reported in Appendix Table~\ref{tab:DWA_regression_aiExecution_GPTsample}.
  The red shaded area marks the 90\% confidence interval around the observed point estimates in the GPT-filtered sample.
  The histograms plot the distribution of average marginal effects across 1{,}000 randomized reshuffles of task positions within occupations.
  Panel~(A) reports results from specification~\eqref{eq:DWA_regression_ai}, while Panels~(B), (C), and (D) augment the baseline specification with SOC major group, SOC minor group, and DWA fixed effects, respectively.
}
\end{figure}




\newpage
\begin{figure}[ht!]
  \begin{center}
  \caption{Effect of Neighboring Tasks' AI Execution Status on Task's AI Automation Likelihood} 
  \label{fig:DWA_regression_aiAutomation_mainSample}
  \begin{subfigure}[b]{\textwidth}
    \captionsetup{labelformat=empty}
    \caption{Panel (A): No Fixed Effects}
    \includegraphics[width=\textwidth]{plots/execTypeVaryingDWA/is_automated/AME_full_is_automated_no_fe_no_dwa.png}
  \end{subfigure}
  
  \vspace{1em}
  
  \begin{subfigure}[b]{\textwidth}
    \captionsetup{labelformat=empty}
    \caption{Panel (B): SOC Major Group Fixed Effects}
    \includegraphics[width=\textwidth]{plots/execTypeVaryingDWA/is_automated/AME_full_is_automated_major_fe_no_dwa.png}
  \end{subfigure}
  
  \vspace{1em}
  
  \begin{subfigure}[b]{\textwidth}
    \captionsetup{labelformat=empty}
    \caption{Panel (C): SOC Minor Group Fixed Effects}
    \includegraphics[width=\textwidth]{plots/execTypeVaryingDWA/is_automated/AME_full_is_automated_minor_fe_no_dwa.png}
  \end{subfigure}

  \vspace{1em}

  \begin{subfigure}[b]{\textwidth}
    \captionsetup{labelformat=empty}
    \caption{Panel (D): Detailed Work Activity Fixed Effects}
    \includegraphics[width=\textwidth]{plots/execTypeVaryingDWA/is_automated/AME_full_is_automated_no_fe_with_dwa.png}
  \end{subfigure}
  \end{center}
  \footnotesize{\emph{Notes:} These graphs show that, among similar tasks appearing in multiple occupations, those whose immediate neighbors are AI-executed exhibit higher probabilities of being automated by AI, while more distant neighbors have little to no effect on a task's AI automation likelihood.
  The red dashed line in each graph indicates the observed average marginal effect of the corresponding variable in the original dataset reported in Table~\ref{tab:DWA_regression_aiAutomation_mainSample}.
  The red shaded area marks the 90\% confidence interval around the observed point estimates in the main sample.
  The histograms plot the distribution of average marginal effects across 1{,}000 randomized reshuffles of task positions within occupations.
  Panel~(A) reports results from specification~\eqref{eq:DWA_regression_ai} but with dependent variable ($\text{is\_automated}_{t}$) instead of ($\text{is\_ai}_{t}$), while Panels~(B), (C), and (D) augment the baseline specification with SOC major group, SOC minor group, and DWA fixed effects, respectively.
}
\end{figure}




\newpage
\begin{figure}[ht!]
  \begin{center}
  \caption{Effect of Neighboring Tasks' AI Execution Status on Task's AI Automation Likelihood (GPT-filtered Sample)} 
  \label{fig:DWA_regression_aiAutomation_GPTsample}
  \begin{subfigure}[b]{\textwidth}
    \captionsetup{labelformat=empty}
    \caption{Panel (A): No Fixed Effects}
    \includegraphics[width=\textwidth]{plots/execTypeVaryingDWA/is_automated/AME_filtered_is_automated_no_fe_no_dwa.png}
  \end{subfigure}
  
  \vspace{1em}
  
  \begin{subfigure}[b]{\textwidth}
    \captionsetup{labelformat=empty}
    \caption{Panel (B): SOC Major Group Fixed Effects}
    \includegraphics[width=\textwidth]{plots/execTypeVaryingDWA/is_automated/AME_filtered_is_automated_major_fe_no_dwa.png}
  \end{subfigure}
  
  \vspace{1em}
  
  \begin{subfigure}[b]{\textwidth}
    \captionsetup{labelformat=empty}
    \caption{Panel (C): SOC Minor Group Fixed Effects}
    \includegraphics[width=\textwidth]{plots/execTypeVaryingDWA/is_automated/AME_filtered_is_automated_minor_fe_no_dwa.png}
  \end{subfigure}

  \vspace{1em}

  \begin{subfigure}[b]{\textwidth}
    \captionsetup{labelformat=empty}
    \caption{Panel (D): Detailed Work Activity Fixed Effects}
    \includegraphics[width=\textwidth]{plots/execTypeVaryingDWA/is_automated/AME_filtered_is_automated_no_fe_with_dwa.png}
  \end{subfigure}
  \end{center}
  \footnotesize{\emph{Notes:} These graphs show that, among similar tasks appearing in multiple occupations, those whose immediate neighbors are AI-executed exhibit higher probabilities of being automated by AI, while more distant neighbors have little to no effect on a task's AI automation likelihood.
  The red dashed line in each graph indicates the observed average marginal effect of the corresponding variable in the original dataset reported in Table~\ref{tab:DWA_regression_aiAutomation_GPTsample}.
  The red shaded area marks the 90\% confidence interval around the observed point estimates in the main sample.
  The histograms plot the distribution of average marginal effects across 1{,}000 randomized reshuffles of task positions within occupations.
  Panel~(A) reports results from specification~\eqref{eq:DWA_regression_ai} but with dependent variable ($\text{is\_automated}_{t}$) instead of ($\text{is\_ai}_{t}$), while Panels~(B), (C), and (D) augment the baseline specification with SOC major group, SOC minor group, and DWA fixed effects, respectively.
}
\end{figure}